% Generated draft: requires mathematical and source review.
\documentclass[11pt]{article}
\usepackage[margin=25mm]{geometry}
\usepackage{fontspec}
\setmainfont{Latin Modern Roman}
\usepackage{amsmath,amssymb,mathtools,unicode-math}
\setmathfont{Latin Modern Math}
\usepackage{enumitem,needspace,xurl,hyperref}
\setlength{\parindent}{0pt}
\setlength{\parskip}{5pt}
\setlength{\emergencystretch}{3em}
\newcommand{\N}{\mathbb{N}}
\newcommand{\Z}{\mathbb{Z}}
\newcommand{\Q}{\mathbb{Q}}
\newcommand{\R}{\mathbb{R}}
\newcommand{\C}{\mathbb{C}}
\newcommand{\F}{\mathbb{F}}
\newcommand{\A}{\mathbb{A}}
\newcommand{\PP}{\mathbb{P}}
\newcommand{\E}{\mathbb{E}}
\newcommand{\eps}{\varepsilon}
\newcommand{\dd}{\,\mathrm{d}}
\begin{document}
\section*{UnsolvedMath ID: 30001288}
\textbf{Source problem code:} catalogue entry
\subsection{Definitions and mathematical statement}
Let \(X\) be a complex algebraic variety, and let \(p\) denote the index occurring in the supplied source notation \(CH_g^p(X)\). Write \(T_{\mathrm W}(X,p)\) solely as a placeholder for the target of Walker's morphic Abel--Jacobi map referred to in that source. The proposed identification is
\[
\operatorname{Alb}\bigl(CH_g^p(X)\bigr)(\C)
\overset{?}{\cong}
T_{\mathrm W}(X,p),
\]
with the isomorphism required to be canonical.

The supplied record presents this as a conjectural identification for a complex algebraic variety. It does not define \(CH_g^p(X)\), the Albanese construction applied to it, or Walker's morphic Abel--Jacobi map and its target. It also does not specify the range of \(p\), whether additional hypotheses on \(X\) are required, or the category and compatibility conditions intended by a canonical isomorphism. Consequently, the displayed question preserves the source notation but remains schematic pending those definitions and qualifications; no smoothness, properness, projectivity, or index restrictions have been added.
\subsection{Short English statement}
The source asks whether the complex points of the Albanese object associated with \(CH_g^p(X)\) are canonically isomorphic to the target of Walker's morphic Abel--Jacobi map. The constructions and their domains require clarification before this can be a fully specified conjecture.
\subsection{Sources}
\begin{itemize}
\item Supplied local UnsolvedMath registry record, permanent ID \(30001288\), exact source code \texttt{catalogue entry}, title: Motivic Albanese and Walker Abel--Jacobi Targets. The record assigns the proposed year 2009 and supplies no named proposer; Walker is named in connection with the map.
\item \emph{Algebraic \(K\)-Theory and Motivic Cohomology} (2009), Oberwolfach Reports, pp.~1731--1774, DOI: \url{https://doi.org/10.4171/owr/2009/31}. The record describes Oberwolfach Reports as documenting mathematical workshops and their open-problem discussions. This report is the supplied authoritative source citation; no precise internal location for this problem is given.
\item Ayoub--Barbieri-Viale on 1-motivic sheaves: \url{https://arxiv.org/abs/math/0607738}. This is the additional supplied literature reference.
\item Literature assessment checked at \texttt{2026-08-21}: the supplied curation labels the problem open, describes the Walker morphic Abel--Jacobi identification as a conjecture, and states that the cited motivic Albanese work does not prove the general identification. This is supplied dated curation derived from the Oberwolfach report, not an independently verified assessment; time-sensitive status and cited claims require independent verification.
\item The registry labels the statement \texttt{corrected\_verified} and reports that its mathematical content was checked against the official Oberwolfach report, with notation, context, and question form normalized without changing the source-level claim. This is supplied verification metadata, not verification performed for this draft.
\end{itemize}
\end{document}
